\lesson{19}{Mar 18 2022 Fri (23:06:17)}{Construct a frequency distribution and a histogram}{Unit 7}

Dr. Martinez is a veterinarian who sees only dogs and cats.
In each appointment, he may or may not give the animal a vaccine.

\begin{center}
  \begin{table}[htpb]
    \centering
    \caption{Two Way Frequency Table that summarized Dr. Martinez's $60$ 
             appointments last week}
    \label{tab:two_way_frequency_table_with_a_random_sample_of_60_items_sold_that_night}
  
    \begin{tabular}{|c|c|c|}
      \hline
      & Dog & Cat \\ 
      \hline
      Vaccine & $13$ & $17$ \\
      \hline
      No vaccine & $21$ & $9$ \\
      \hline
    \end{tabular}
  \end{table}
\end{center}

Let no vaccine be the event that a randomly chosen appointment (from the table
\ref{tab:two_way_frequency_table_with_a_random_sample_of_60_items_sold_that_night})
did not include a vaccine.

Let cat be the event that a randomly chosen appointment (from the table
\ref{tab:two_way_frequency_table_with_a_random_sample_of_60_items_sold_that_night})
involved a cat.


Let's try and find the following properties:

\begin{itemize}
  \label{item:two_way_frequency_table}

  \item \textbf{P(no vaccine)} = \boxed{ }

    This is the probability that a randomly chosen appointment (from the table
    \ref{tab:two_way_frequency_table_with_a_random_sample_of_60_items_sold_that_night})
    did not include a vaccine.

    We must divide the number of appointments that did not include a vaccine by
    the total number of appointments ($60$).

    So, here's our computation:

    \begin{equation}
      \label{eqt:no_vaccine}
      \begin{split}
        P(\textrm{no vaccine}) &= \frac{\textrm{number of appointments with no
          vaccine}}{\textrm{total number of appointments}} \\
                               &= \frac{21 + 9}{60} \\
                               &= \frac{30}{60} \\
                               &= boxed{0.5} \\
      \end{split}
    .\end{equation}

    We get \textbf{P(no vaccine)} $= 0.5$.

  \item \textbf{P(cat and no vaccine)} = \boxed{ }

    This is the probability that a randomly chosen appointment (from the table
    \ref{tab:two_way_frequency_table_with_a_random_sample_of_60_items_sold_that_night})
    did not include a vaccine and involved a cat.

    We must divide the number of appointments that did not include a vaccine and
    involved a cat by the total number of appointments ($60$).

    So, here's our computation:
    
    \begin{equation}
      \label{eqt:cat_and_no_vaccine}
      \begin{split}
        P( \textrm{cat and no vaccine}) &= \frac{\textrm{number of cat
          appointments with no vaccine}}{\textrm{total number of appointments}} \\
                                        &= \frac{9}{60} \\
                                        &= \boxed{0.15} \\
      \end{split}
    .\end{equation}

    We get \textbf{P(cat and no vaccine)} $= 0.15$.

  \item \textbf{P(cat | no vaccine)} = \boxed{ }

    This is the probability that a randomly chosen appointment (from the table
    \ref{tab:two_way_frequency_table_with_a_random_sample_of_60_items_sold_that_night})
    involved a cat given that it didn't include a vaccine.

    So, we look only at appointments that didn't include a vaccine. There are
    $21 + 9 = 30$ such appointments.

    So, here's our computation:

    \begin{equation}
      \label{eqt:cat_no_vaccine}
      \begin{split}
        P(\textrm{cat | no vaccine}) &= \frac{\textrm{Number of cat appointments
          with no vaccine}}{\textrm{total number of appointments with no
          vaccine}} \\
                                     &= \frac{9}{30} \\
                                     &= \boxed{0.3} \\
      \end{split}
    .\end{equation}

    we get \textbf{P(cat | no vaccine)} $= 0.3$.
\end{itemize}

\newpage
